\documentclass[11pt]{article}
\usepackage{iftex}
\ifPDFTeX
\usepackage[T1]{fontenc}
\usepackage[utf8]{inputenc}
\else
\usepackage{fontspec}
\fi
\usepackage{amsmath,amssymb,amsthm}
\usepackage[margin=1in]{geometry}
\usepackage[hidelinks]{hyperref}
\newtheorem{theorem}{Theorem}
\newtheorem{lemma}[theorem]{Lemma}
\newtheorem{corollary}[theorem]{Corollary}
\newtheorem{proposition}[theorem]{Proposition}
\newcommand{\C}{\mathcal C}
\newcommand{\G}{\mathcal G}
\newcommand{\MA}{\operatorname{MA}}
\newcommand{\DD}{d'd''}
\newcommand{\R}{\mathbb R}
\newcommand{\N}{\mathbb N}
\setlength{\parindent}{0pt}
\setlength{\parskip}{0.5em}
\emergencystretch=3em
\title{Quadratic energy scaling, solid full duals, and variational coercivity\\on the actual nonalgebraic analytic surface}
\author{Unreviewed scoped progress on original P16.7}
\date{9 October 2026}
\begin{document}
\maketitle
\begin{abstract}
The uniform reference-skeleton bound for normalized continuous weights
turns the exact cubic energy formula into a quadratic estimate at small
positive scaling. We deduce an equivalent bounded-unit-energy test for
the full Legendre dual, square-root bounds on every normalized energy
sublevel, and coercivity of every finite-full-dual variational functional.
The class of positive mass-two finite-full-dual Radon measures is solid
under bounded domination. In particular every total-dimension-two point
supports a normalized Dirac measure of finite full dual energy, using
an actual bounded Monge--Ampere atom. Countable mixtures with an explicit
summable dual cost inherit finite dual energy and uniform integrability.
These results use the qualified actual continuous calculus and Gauss/atom
foundations; they do not establish nonordered analytic compactness,
envelopes, or prescribed-measure existence. Original P16.7 remains open.
\end{abstract}

\section{Actual object and precise analytic foundations}
Use the actual compact analytic quotient
\[
 K=\mathbb C((s)),\quad |s|=e^{-1/8},\quad t=s^8,
 \qquad X=(\mathbb G_m^2)^{\mathrm{an}}_K/
                    \langle(t,2),(3,t)\rangle.
\]
On injective torus charts write $r_i=-\log|z_i|$,
$\rho=|r|^2/2$, and $\theta=\DD\rho$. Reference changes are
affine unit logarithms. The raw positive mass is $\int_X\theta^2=2$.
Let $\C$ be the entire continuous $\theta$-psh class obtained from the
minimal continuous CLD psh sheaf, including all its continuous stages.
For $a\in\C$, put $\theta_a=\theta+\DD a$,
$\MA(a)=\theta_a^2$, and
\[
 E(a)=\frac13\int_X a[\theta^2+\theta\wedge\theta_a+\theta_a^2].
\]
The previously reconstructed continuous Green calculus supplies the
positive bilinear gradient pairing $Q_T$ for positive closed $(1,1)$
currents $T$, and the identities
\begin{equation}\label{green}
 Q_T(f,g)=-\int_X f\,\DD(gT)=-\int_X g\,\DD(fT),
 \qquad Q_T(f,f)\ge0.
\end{equation}
These identities concern the actual continuous strong-gradient domain,
not a presumed pullback of an arbitrary current. Constants and convex
combinations preserve $\C$; all bounded products have mass two.
The resulting energy is monotone and concave, with
\begin{equation}\label{translation}
 E(0)=0,\qquad E(a+c)=E(a)+2c.
\end{equation}
For the segment $ta$, $0\le t\le1$, first variation and Green give
\begin{equation}\label{variation}
 \frac{d}{dt}E(ta)=\int_X a\,\MA(ta),\qquad
 \frac{d^2}{dt^2}E(ta)=-2Q_{\theta_{ta}}(a,a).
\end{equation}
The whole continuous calculus and actual quotient are qualified analytic
foundations of this application. They are reconstructed in current
working papers 402/403/409, and are not inferred from a compiler,
historical verdict, or source hash. The new arguments below are complete
given \eqref{green}--\eqref{variation} and the two geometric facts stated
next. The full recursive audit of all deeper analytic foundations is
not newly claimed here.

\subsection{The reference bound used in the scaling argument}
Let $S$ be the joint coordinate Gauss skeleton. For every $a\in\C$,
the actual annulus and polyannulus comparisons give
\begin{equation}\label{gauss}
 \sup_Xa=\max_S a,\qquad
 \sup_Xa=0\Longrightarrow -\tfrac12\le a|_S\le0,
 \qquad \theta^2=\sigma_*(2\,dr_1dr_2).
\end{equation}
Here $\sigma$ is the Gauss section on $\R^2/\mathbb Z^2$.
We recall the precise application rather than assuming an analytic map
obtained by varying a higher-dimensional Shilov section. Fix the second
Gauss coordinate and take the actual first-coordinate analytic curve
over its completed residue field. The pullback of $\rho+a$ is
classically subharmonic. Comparison on any closed annulus with the
affine harmonic interpolation of its two endpoint Gauss values shows
that $q^2/2+a(\sigma(q,r_2))$ is convex. The deck translations give
period one for $a(\sigma(q,r_2))$. Secants to $q-1$ and $q+1$ bound
every convex subgradient between $q-1/2$ and $q+1/2$. Thus this periodic
function is $1/2$-Lipschitz. The other actual coordinate curve gives
the same assertion in the second coordinate. The sum of the two maximal
circle distances bounds the oscillation on $S$ by $1/2$.
The maximum principle on successively shrinking rectangular polyannuli,
whose four Shilov points are the joint Gauss vertices, gives
$a(x)\le a(\sigma(r(x)))$ and hence the first two assertions of
\eqref{gauss}.

For the measure assertion, a point off the coordinate skeleton has
strict cancellation in some Laurent polynomial. On a neighborhood
where that strict inequality persists, its coordinate tropicalization
lies in the polynomial's finite tropical hypersurface of dimension at
most one. The coordinate $(2,2)$ form vanishes there by the actual
moment-form dimension rule. Thus $\theta^2$ is supported on $S$.
The identity moment has calibration one; its Hessian is the identity
and the raw wedge factor is $2!$. The actual measure is therefore
twice Lebesgue measure on the fundamental square. The full immediate
CLD integration/calibration and classical GR comparison applications
are qualified as in working paper407. None is extended to a varying
Shilov section morphism. Consequently
\begin{equation}\label{A}
 A(a):=\int_X(-a)\theta^2\le1
 \quad\text{if }a\in\C,\ \sup_Xa=0.
\end{equation}

\subsection{Actual point atoms and generated weights}
At each point $x$ of total graded dimension two the qualified actual
point construction supplies $w_x\in\C$ and $0<\alpha_x\le2$ such that
\begin{equation}\label{atom}
 \nu_x:=\MA(w_x)\ge\alpha_x\delta_x,
 \qquad E^*(\nu_x)<\infty.
\end{equation}
This is the whole actual bounded measure, not only its tropical image.
Its construction selects two algebraically independent homogeneous
Laurent initial forms at $x$. Their map sends $x$ to a weighted Gauss
point and has a finite flat local representative with singleton fibre.
The raw top mass of $\max(0,\log(|P_1|/r_1),\log(|P_2|/r_2))$ is one
at the target vertex: the gradient simplex has area $1/2$, and the
wedge factor is $2!$. Finite properness and the singleton fibre turn
the smooth approximating measures into an actual atom of local degree
$D>0$ at $x$. A quadratic-reference cap, zero on an outer collar,
globalizes the weight, with $\alpha_x=\epsilon^2D>0$. The complete
initial-form, finite-flat, calibration, and all-stage cap proof is the
qualified actual construction in working paper409. No ground-field
extension or fraction-field valuation through nonzero support is
substituted for this construction. Bounded tangent gives the full dual
statement in \eqref{atom}. Finite dual for $2\delta_x$ will be a new
conclusion below, not an assumption of the atom construction.

The original generated class $\G$ consists of actual pointwise
$a_j\downarrow u$, $a_j\in\C$, $\inf_jE(a_j)>-\infty$.
Its actual energy and measure are presentation independent under the
qualified current calculus. Every positive mass-two finite-full-dual
measure $\mu$ satisfies
\begin{equation}\label{generatedD}
 u\le0\Longrightarrow
 \int_X(-u)\,d\mu\le E^*(\mu)-E(u).
\end{equation}
For completeness, the same dual inequality first holds for the bounded
presentation terms. Subtract a common continuous upper bound and use
ordinary sequential monotone convergence. This proves both integrability
and \eqref{generatedD}; generated pole sets are $\mu$-null. This is a
sequential argument and invokes no monotone convergence theorem for
arbitrary measurable nets.

\section{Exact cubic identity and quadratic small-scale control}
\begin{lemma}\label{scale}
For $a\in\C$ with $\sup_Xa=0$, set
$C=-E(a)$, $q_0=Q_\theta(a,a)$, and $q_1=Q_{\theta_a}(a,a)$.
Then $A=A(a),q_0,q_1,C$ are nonnegative, and
\begin{align}
 C&=A+\frac23q_0+\frac13q_1,\label{coeff}\\
 -E(ta)&=tA+t^2(1-t/3)q_0+\frac{t^3}{3}q_1
       \le t+\frac32t^2 C\quad(0\le t\le1).\label{small}
\end{align}
\end{lemma}
\begin{proof}
The pairing is linear in the current and
$\theta_{ta}=(1-t)\theta+t\theta_a$. Therefore
\eqref{variation} gives $E''(ta)=-2((1-t)q_0+tq_1)$,
where the derivative is with respect to $t$. At zero the value is zero
and the first derivative is $\int a\theta^2=-A$.
Integrating this affine second derivative twice gives exactly
$-E(ta)=tA+t^2q_0-t^3q_0/3+t^3q_1/3$.
Putting $t=1$ proves \eqref{coeff}. All three coefficients there
are nonnegative by positivity and $a\le0$.
For $0\le t\le1$, the last two terms are at most
$t^2(q_0+q_1/3)\le(3/2)t^2(C-A)$.
Use $A\le1$ and $C-A\le C$ to obtain \eqref{small}.
No regularity beyond the actual bounded Green calculus is used.
\end{proof}

\section{An equivalent bounded-unit-energy criterion for the full dual}
For a positive Radon measure define the full bounded Legendre dual by
\[
 E^*(\mu)=\sup_{a\in\C}\bigl(E(a)-\int_Xa\,d\mu\bigr).
\]
Constants in \eqref{translation} show that finite dual forces
$\mu(X)=2$. For such a measure the functional is unchanged by adding
a constant, so the supremum may be taken over $\sup_Xa=0$.
Define the bounded-unit-energy test
\[
 N(\mu)=\sup\left\{\int_X(-a)\,d\mu:
       a\in\C,\ \sup_Xa=0,\ E(a)\ge-1\right\}.
\]
\begin{theorem}\label{unit}
For every positive Radon measure of mass two,
$E^*(\mu)<\infty$ if and only if $N(\mu)<\infty$.
When finite, the quantitative bounds are
\begin{equation}\label{unitbound}
 N(\mu)\le E^*(\mu)+1,
 \qquad E^*(\mu)\le2N(\mu)+\frac32N(\mu)^2.
\end{equation}
The tests are the whole original continuous class, with arbitrary
measure support.
\end{theorem}
\begin{proof}
The first inequality is the defining full dual inequality.
For the converse fix normalized $a$ and $C=-E(a)\ge0$. Put
\[
 \tau=\frac1{2\max(1,\sqrt{3C/2})}.
\]
Then $0<\tau\le1/2$, $\tau a\in\C$ is normalized, and
\eqref{small} gives $-E(\tau a)\le1/2+1/4=3/4$.
Thus $\tau\int(-a)d\mu\le N(\mu)$ and
\[
 \int(-a)d\mu\le 2N(\mu)\max(1,\sqrt{3C/2})
              \le 2N(\mu)+\sqrt6 N(\mu)\sqrt C.
\]
Subtract $C$ and maximize over $C\ge0$:
$\sup_{C\ge0}(-C+\sqrt6 N\sqrt C)=3N^2/2$.
Constant invariance then proves the full second bound in
\eqref{unitbound}, not merely a normalized or finite set of tests.
\end{proof}

\section{Square-root integrability and universal variational coercivity}
\begin{theorem}\label{root}
Let $\nu$ be any positive mass-two Radon measure with
$B=E^*(\nu)<\infty$. Then $B\ge0$ and for every normalized
$a\in\C$, $C=-E(a)$,
\begin{equation}\label{rootbound}
 \int_X(-a)\,d\nu\le1+B+\sqrt{6BC}.
\end{equation}
In particular $F_\nu(a)=E(a)-\int a\,d\nu$ satisfies
\begin{equation}\label{coercive}
 F_\nu(a)\le-C+1+B+\sqrt{6BC}\longrightarrow-\infty
 \quad\text{as }C\longrightarrow\infty.
\end{equation}
If $F_\nu(a)\ge-\eta$ with $\eta\ge0$, then
\begin{equation}\label{sublevel}
 -E(a)\le\frac14\left(\sqrt{6B}+
                       \sqrt{10B+4+4\eta}\right)^2.
\end{equation}
\end{theorem}
\begin{proof}
The zero competitor gives $B\ge0$. Apply the full dual inequality
to $ta$ and use the stronger intermediate form of \eqref{small}:
\[
 \int(-a)d\nu\le A+B/t+(3/2)tC \quad(0<t\le1).
\]
If $B=0$, let $t\downarrow0$ to get the stronger bound
$\int(-a)d\nu\le A\le1$. If $C=0$, take $t=1$.
For $B,C>0$ with $B<(3/2)C$, choose
$t=\sqrt{2B/(3C)}<1$; the two variable terms sum to
$\sqrt{6BC}$. If $B\ge(3/2)C$, choose $t=1$;
then $(3/2)C\le\sqrt{6BC}$. These cases prove
\eqref{rootbound}. Subtracting $C$ proves \eqref{coercive}.
For \eqref{sublevel}, put $y=\sqrt C$. The assumed lower
functional bound gives $y^2-\sqrt{6B}y-(1+B+\eta)\le0$.
Its nonnegative root is the bound displayed in \eqref{sublevel}.
\end{proof}

Since $F_\nu(0)=0$ and $\sup F_\nu=B$, a bounded-class maximizing
sequence exists by the definition of supremum and may be normalized by
constants. Its tail has $F_\nu\ge-1$, so \eqref{sublevel} with
$\eta=1$ bounds its energies. This establishes coercivity for EVERY
finite-full-dual measure. It does not establish an actual convergent
psh representative or upper semicontinuity of the objective on a
nonordered analytic completion.

The same bounds hold for normalized $u\in\G$ by normalizing decreasing
continuous presentations. Indeed, for an arbitrary presentation
$a_j\downarrow u$, the real maxima $s_j=\sup_Xa_j$ decrease to
$\sup_Xu$. To verify this assertion, choose maximum points of $a_j$;
compactness gives a cluster subnet, and each fixed earlier continuous
term bounds every late maximum. At the cluster point the infimum is
at least $\inf s_j$, giving the reverse inequality. For $\sup u=0$
we have $s_j\downarrow0$ and normalized $b_j=a_j-s_j\to u$
pointwise; $E(b_j)\to E(u)$. Each $\int(-b_j)d\nu$ converges to
$\int(-u)d\nu$ by the previously established sequential $L^1$
convergence of $a_j$ and $s_j\to0$. The $b_j$ need not decrease,
and no such assertion is required. Passing \eqref{rootbound} gives
the claimed extension. The coercivity and energy bound extend too.

\section{Solidity under bounded domination and full-dimensional Dirac tests}
\begin{theorem}\label{solid}
Suppose $\mu,\nu$ are positive mass-two Radon measures,
$\mu\le\lambda\nu$, and $B=E^*(\nu)<\infty$.
Then $\lambda\ge1$ and
\begin{equation}\label{solidbound}
 E^*(\mu)\le\lambda(1+B)+\frac32\lambda^2B<\infty.
\end{equation}
\end{theorem}
\begin{proof}
Masses give $\lambda\ge1$. For normalized $a\le0$ with
$C=-E(a)$, positivity and \eqref{rootbound} give
\[
 E(a)-\int a\,d\mu
 \le-C+\lambda(1+B)+\lambda\sqrt{6BC}.
\]
The supremum of the last variable expression is $3\lambda^2B/2$.
Constants preserve the full dual objective, proving
\eqref{solidbound}. Merely scaling a component of $\nu$ and using
the original affine dual bound would not prove this: the square-root
growth just established is the essential extra step.
\end{proof}

\begin{corollary}\label{dirac}
Every point $x\in X$ of total graded dimension two has
$E^*(2\delta_x)<\infty$. With the actual atom in \eqref{atom},
write $B_x=E^*(\nu_x)$ and $\lambda_x=2/\alpha_x$.
An explicit bound is
\begin{equation}\label{Lx}
 E^*(2\delta_x)\le L_x:=
             \lambda_x(1+B_x)+\frac32\lambda_x^2B_x.
\end{equation}
Every normalized $u\in\G$ of energy at least $-C$ therefore satisfies
$0\le-u(x)\le(L_x+C)/2$.
\end{corollary}
\begin{proof}
The actual measure inequality gives $2\delta_x\le\lambda_x\nu_x$;
both measures have raw mass two. Apply Theorem~\ref{solid}.
Then use \eqref{generatedD} for $2\delta_x$. This conclusion is
stronger than finiteness of each individual generated value obtained
from the unnormalized atom alone. It uses \eqref{gauss} and the
scaling theorem and is not a premise of paper409.
\end{proof}
This asserts finite dual energy, not the existence of a weight whose
entire Monge--Ampere measure is $2\delta_x$.

\section{Countable mixtures and a precise uniform-integrability sector}
\begin{proposition}\label{mixture}
Let $\mu_n$ be positive mass-two Radon measures,
$p_n\ge0$, $\sum_n p_n=1$, and
$B_n=E^*(\mu_n)<\infty$. If $\sum_n p_nB_n<\infty$, then
$\nu=\sum_n p_n\mu_n$ is an actual positive Radon measure of mass
two and
\begin{equation}\label{mixdual}
 E^*(\nu)\le\sum_n p_nB_n.
\end{equation}
Suppose additionally that each $\mu_n$ has uniform negative-part tails
on the normalized energy sublevel $\{u\in\G:\sup u=0,E(u)\ge-C\}$.
Then so does $\nu$, and every mass-two $\mu\le\lambda\nu$ has finite
full dual and the same uniform-tail property.
\end{proposition}
\begin{proof}
The sum converges in total variation since all masses are two and
$\sum p_n=1$. Radon inner and outer regularity follows by first
approximating a finite initial sum and then using the total-variation
bound on its tail; thus the sum is Radon. For a bounded test $a$,
\[
 E(a)-\int a\,d\nu
   =\sum_n p_n\left(E(a)-\int a\,d\mu_n\right)
   \le\sum_n p_nB_n.
\]
The equality holds by absolute convergence of integrals of the bounded
test, and the right side is finite. Take the supremum over all $a$.
For normalized generated $u$ in the indicated sublevel,
\eqref{generatedD} gives $\int(-u)d\mu_n\le B_n+C$.
Consequently the contribution of all $n>N$ to ANY negative-part tail
is at most $\sum_{n>N}p_n(B_n+C)$, independently of $u$ and the
tail threshold. This tends to zero. The finitely many initial tails
vanish uniformly by assumption, proving the assertion for $\nu$.
All exchanges here are countable nonnegative sums, justified by
ordinary Tonelli; there is no measurable-net limit. Finally bounded
domination gives the tail estimate for $\mu$, and Theorem~\ref{solid}
gives its finite full dual.
\end{proof}

Each fixed bounded $\MA(w)$ has the qualified uniform-tail theorem
of paper407. Thus Proposition~\ref{mixture} applies to countable
mixtures of such measures with summable full-dual costs and to their
boundedly dominated mass-two measures. This is a sufficient sector
with arbitrary support allowed by the actual tests, not a statement
that every finite-full-dual measure admits this domination.

\begin{corollary}\label{countable}
For any specified countable sequence of total-dimension-two points
$x_n$, there are strictly positive weights $p_n$, summing to one, such
that $\nu=2\sum_n p_n\delta_{x_n}$ has finite full dual and uniform
negative-part tails on every normalized finite energy sublevel.
\end{corollary}
\begin{proof}
Choose the finite bounds $L_{x_n}$ in \eqref{Lx} and normalize
$p_n=c\,2^{-n}/(1+L_{x_n})$, $n\ge1$, by the positive finite sum.
Then $\sum p_nL_{x_n}<\infty$. Each point test has a uniform
sublevel value bound $(L_{x_n}+C)/2$, so its negative-part tail is
identically zero beyond that bound. Apply Proposition~\ref{mixture}.
No density or countability property of $X$ is assumed; this constructs
a measure on the specified sequence and does not establish a
full-support countably atomic measure.
\end{proof}

\section{Scope and unchanged original residual}
The new algebraic identities and inequalities establish the full
bounded-unit criterion, universal square-root dual integrability,
variational coercivity, solid domination, actual point-Dirac finite
duals and the stated summable-mixture sector. They retain raw mass two,
all continuous stages and arbitrary measure support. The generated
extensions use actual sequential presentations and ordinary $L^1$
convergence. No prescribed measure is restricted to an existing
Monge--Ampere image.

The actual quotient, entire strong-gradient/continuous product calculus,
GR recognition and annulus/Shilov comparison, actual CLD calibrated
measure and local point-atom constructions remain qualified analytic
foundations. Their current applications were checked against the whole
relevant current407/409 bodies. Source identity and repeated retention
are not extra independent source readings. No newly completed recursive
audit of all deeper analytic/algebraic sources is asserted.

The nonordered energy-sublevel representative and compactness theorem
remains unresolved. So do adequate intrinsic finite-energy recognition,
actual envelopes/orthogonality or sufficient replacements, and existence
for EVERY prescribed positive mass-two finite-full-dual Radon measure.
Coercivity confines maximizing sequences to a sublevel; it supplies no
actual limit, semicontinuity, differentiable envelope variation or
maximizer. Uniform integrability was proved only in the explicit
summable-mixture/dominated sector, not for every finite-full-dual
measure. No original Dirac equation is solved here. No metrizability,
countable dense set, minimum closure, positive USC regularization,
arbitrary measurable-net monotone convergence or full-support
countable Radon measure is assumed. This complete scoped working paper
is UNREVIEWED; original P16.7 remains OPEN and the solved count is
unchanged.
\end{document}
